% !Mode:: "TeX:UTF-8"

\chapter{系统分析}\label{ch:3}

\section{可行性分析}

\subsection{技术可行性}

从实现条件看，本系统具备可行性。后端使用的 SpringBoot 3.3、MyBatis-Plus、SpringSecurity 和 Spring AI 都属于 Spring 生态中的主流组件，版本成熟，文档完善，社区活跃。前端使用的 Vue3、Element Plus、ECharts 和 Axios 也已有成熟的应用基础。AI 能力通过 Spring AI 接入外部模型服务，不依赖特定厂商 SDK，模型切换成本低。开发环境方面，普通配置的办公电脑即可同时运行前后端服务和数据库，硬件门槛不高。

\subsection{经济可行性分析}

从成本角度看，本系统的开发投入集中在开发时间和基础设备上。技术栈均基于开源框架，无需购买商业许可。AI 模型调用可采用兼容方式接入免费或低成本模型服务，初期测试和运行成本较容易控制。学校现有服务器环境即可完成部署，不需要额外购置软硬件。系统投入使用后，体测数据录入、统计和查询过程中的重复性人工操作会减少，管理人力成本也会下降。整体上具备经济可行性。

\subsection{操作可行性分析}

从操作层面看，本系统面向管理员、教师和学生三类角色，各自的业务边界比较清楚。管理员负责用户、课程和全局数据维护。教师负责课程内体测录入、修正和审批。学生负责个人查询和请假申请。系统界面围绕既有教学管理流程展开，用户不需要改变原有业务顺序就能完成主要操作，在高校体育教学场景中具有操作可行性。

\section{功能需求分析}

本系统面向高校体测管理业务，服务管理员、教师和学生三类用户。管理员负责基础数据和全局配置，教师负责本人课程范围内的体测录入、成绩复核和请假审批，学生负责个人查询和请假申请。围绕三类用户，系统需形成从课程建立、数据采集、成绩计算到结果查询和分析追溯的完整流程。

（1）管理员功能需求。管理员需要完成教师账号、角色权限、课程信息、学生基础信息和日志审计等全局管理工作。系统应支持教师账号新增、修改、禁用和密码重置，支持角色与菜单权限配置。管理员可通过 Excel 批量导入教师、学生和课程数据，也可按学期、课程或班级导出体测成绩。课程管理方面，管理员需为课程分配教师并维护课程状态，为后续体测录入和成绩计算提供基础。

（2）教师功能需求。教师登录后只能查看和操作本人负责的课程。教师可按课程、班级和体测项目录入或批量导入成绩，支持异常数据修正、备注补充和模板导出。课程结束时，系统需根据已录入体测记录计算总分、等级和达标情况。课程结束后如需补录，应先恢复课程状态再重新计算成绩。教师还需审批本人课程下学生提交的请假申请，审批过程记录审批人、审批时间、审批状态和审批意见。

（3）学生功能需求。学生通过独立入口使用学号和密码登录，首次登录后需要修改默认密码。登录后，学生只能查看与本人学号相关的数据，包括历史体测成绩、课程与体测安排、个人健康建议和请假记录。请假申请中课程信息与本人课程绑定，提交后可查看审批进度与结果。

（4）核心业务流程需求。系统需支持“课程创建与分配—体测模板导出—成绩录入或导入—数据校验与修正—课程结束算分—统计分析—学生查询”的主流程，同时支持“学生提交请假—教师按课程审批—特殊标签写入体测备注—审批留痕”的辅助流程。对兵役、免测等特殊申请，审批通过后应在对应体测记录中形成备注，保证后续成绩核对有据可查。

（5）AI 分析与问答需求。管理员和教师可查看指定学期或课程的达标率、等级分布和项目表现，并发起 AI 分析。AI 问答需支持按学生、课程、班级和学期等范围组织证据，对证据不足的问题应明确提示。分析和问答结果写入日志或分析记录，便于复核。

（6）数据管理与追溯需求。系统应记录账号维护、权限调整、导入导出、体测修改、课程结束、请假审批和 AI 分析等关键操作，支持按人员、时间和类型进行查询。体测数据、成绩结果、请假审批和 AI 输出之间保持可追踪关系。

\section{系统的非功能需求}

性能和运行目标方面，常规校园网络环境下常用页面应保持流畅的交互体验，登录、列表查询、体测历史读取和问答结果返回不应长时间阻塞。系统需支持按课程或班级进行批量导入导出，并在导入过程中完成格式、范围和唯一性校验，以适应本校普通体测与集体体测两类组织方式下的数据处理需求。

安全性需求方面，系统采用登录认证与角色授权相结合的方式控制访问范围。管理员拥有全局管理权限，教师只能操作本人负责的课程、体测记录和请假申请，学生只能查看本人数据并提交个人申请。密码采用安全散列方式存储，接口层需校验用户身份、角色权限和数据归属，防止越权查询或审批。

可靠性需求方面，系统需对手工录入和批量导入的数据进行完整性、格式和范围校验。课程结束、成绩计算、请假审批和特殊标签写回等操作应保证数据状态一致。发生异常时，系统应给出明确提示，并通过操作日志、审批记录和 AI 分析日志定位问题。

可维护性需求方面，系统应保持前端页面、后端接口、业务服务和数据库表结构之间的职责清晰。AI 模块应通过统一服务入口接入模型调用，并记录调用异常和回退结果，便于后续排查和维护。

运行环境与兼容性需求方面，服务器端可运行于常见的 Java 应用部署环境，客户端需兼容 Chrome、Edge 等主流浏览器，保证管理员端、教师端和学生端页面均能正常访问。

\section{主要业务流程图}

业务流程从课程准备开始。管理员完成课程与教师分配后，教师登录系统查看负责班级并录入或导入体测数据，课程结束后触发成绩计算和结果归档。流程如图 \ref{fig:business-flow} 所示。

\begin{figure}[H]
  \centering
  \includegraphics[width=0.9\textwidth]{fig3_1_business_flow.png}
  \caption{业务流程图}
  \label{fig:business-flow}
\end{figure}

\section{系统用例图}

管理员、教师和学生分别围绕各自角色职责参与系统业务。管理员主要负责用户、课程、成绩和日志等全局功能，教师主要负责课程体测数据与审批业务，学生主要负责个人查询和请假申请。三类角色的用例如图 \ref{fig:admin-usecase}、图 \ref{fig:teacher-usecase} 和图 \ref{fig:student-usecase} 所示。

\begin{figure}[H]
  \centering
  \includegraphics[width=0.72\textwidth]{fig3_2_usecase_admin.png}
  \caption{管理员用例图}
  \label{fig:admin-usecase}
\end{figure}

\begin{figure}[H]
  \centering
  \includegraphics[width=0.72\textwidth]{fig3_3_usecase_teacher.png}
  \caption{教师用例图}
  \label{fig:teacher-usecase}
\end{figure}

\begin{figure}[H]
  \centering
  \includegraphics[width=0.72\textwidth]{fig3_4_usecase_student.png}
  \caption{学生用例图}
  \label{fig:student-usecase}
\end{figure}

\section{数据流图}

系统数据流围绕课程、学生、体测记录、成绩统计、请假审批和 AI 分析展开。顶层数据流体现外部参与者与系统之间的信息交换，第一层数据流展示各功能模块之间的数据处理关系，如图 \ref{fig:context-data-flow} 和图 \ref{fig:level1-data-flow} 所示。

\begin{figure}[H]
  \centering
  \includegraphics[width=0.9\textwidth]{fig3_5_context_data_flow.png}
  \caption{系统顶层数据流图}
  \label{fig:context-data-flow}
\end{figure}

\begin{figure}[H]
  \centering
  \includegraphics[width=0.9\textwidth]{fig3_6_level1_data_flow.png}
  \caption{系统第一层数据流图}
  \label{fig:level1-data-flow}
\end{figure}
